\fullwidth{\Large \textbf{Determinantes e Álgebra Linear no $\R^n$}}


\begin{pycode}
import sympy as sp

def pt(Y):
 return((Y[0],Y[1],Y[2]))

A=sp.Matrix([3,2,0]).T
B=sp.Matrix([1,0,2]).T
C=sp.Matrix([0,4,3]).T

AB=B-A
AC=C-A
u=AB.cross(AC)
n=u.norm()
area=n/2

\end{pycode}


\question Determine a área do triângulo cujos vértices são os pontos $A=\pyl{pt(A)}$, $B=\pyl{pt(B)}$ e $C=\pyl{pt(C)}$.

\begin{solution}
A área do triângulo é dada por
\[\frac{1}{2}\|\vt{AB}\times \vt{AC}\|.\]
Como $\vt{AB}=\pyl{pt(AB)}$ e $\vt{AC}=\pyl{pt(AC)}$, então 
$\vt{AB}\times \vt{AC}=\pyl{pt(u)}$ e portanto a área do triângulo é dada por
\[\frac{1}{2}\|\pyl{pt(u)}\|=\pyl{area}.\]

\end{solution}

\begin{pycode}
import sympy as sp

def pt(Y):
 return((Y[0],Y[1],Y[2]))

x,y,z,t=sp.symbols('x y z t',real=True)

X=sp.Matrix([x,y,z])

n1=sp.Matrix([1,1,1])
p1=sp.Eq(n1.dot(X),2)

n2=sp.Matrix([1,-1,-1])
p2=sp.Eq(n2.dot(X),0)

d=sp.sqrt(6)

vs=sp.Matrix([1,-1,1])
ps=sp.Matrix([1,0,-1])

B=sp.Matrix([2,0])

Aa=sp.Matrix.hstack(sp.Matrix.vstack(n1.T,n2.T),B)
R=Aa.rref()[0]

A=sp.Matrix([1,1-t,t])

AP=ps-A
u=AP.cross(vs)
nv=vs.norm()

sol=sp.solve(sp.Eq(sp.expand(u.dot(u)),(nv*d)**2),t)
def aum(M, m):
 # Cria a parte do array baseada no número de colunas
 # As primeiras (m-1) colunas são normais, a última é separada por |
 col = 'c'*(m-1)+'|c'
 MM = sp.latex(M).replace(r'\begin{matrix}', r'\begin{array}{'+col+ '}')
 MM = MM.replace(r'\end{matrix}', r'\end{array}')
 return MM
\end{pycode}
%-------------------------------------------------------

\question 
\begin{parts}
\part Obtenha a reta $r$ dada pela interseção dos planos $\pi_1: \pyl{p1}$ e $\pi_2: \pyl{p2}$.

\part Determine os pontos da reta $r$ que estão à distância $\pyl{d}$ da reta 
$s:\begin{cases}
x=\pyl{ps[0]+vs[0]*t}\\
y=\pyl{ps[1]+vs[1]*t}\\
z=\pyl{ps[2]+vs[2]*t}, \ t\in \R.
\end{cases}$
\end{parts}


\begin{solution}
\begin{parts}
\part Basta resolvermos os sistema
\[\begin{cases}
\pyl{p1}\\ \pyl{p2}.
\end{cases}\]
Escalonando à forma reduzida:
\[
\py{aum(Aa,4)}\sim \py{aum(R,4)} \Rightarrow \begin{cases}
x=1\\ y=1-z
\end{cases}
\]
Fazendo $z=t$, obtemos 
\[r:\begin{cases}
x=1\\ y=1-t\\ z=t,\ t\in \R.
\end{cases}\]

\part Portanto, um ponto genérico da reta $r$ é da forma $A=\pyl{pt(A)}$. Queremos determinar $t$ tal que $d(A,s)=\pyl{d}$. Sabemos que $P=\pyl{pt(ps)}\in s$ e que $\vec{v}=\pyl{pt(vs)}$ é um vetor diretor de $s$. Com isso, 
\[d(A,s)=\frac{\|\vt{AP}\times \vec{v}\|}{\|\vec{v}\|}.\]
Como $\vt{AP}=\pyl{pt(AP)}$, então $\vt{AP}\times \vec{v}=\pyl{pt(u)}$ e portanto
\[\pyl{d}=\frac{\|\pyl{pt(u)}\|}{\pyl{nv}}=\frac{\pyl{u.norm()}}{\pyl{nv}}.\]
Donde,
\[\pyl{sp.expand(u.dot(u))}=\pyl{(nv*d)**2}\Rightarrow t=\pyl{sol[0]} 
\text{ ou } t=\pyl{sol[1]}.\]
Logo, os pontos buscados são:
\[A_1=\pyl{pt(A.subs(t,sol[0]))} \text{ e } A_2=\pyl{pt(A.subs(t,sol[1]))}.\]


\end{parts}
\end{solution}

%----------------------------------------------------------------
\question 	Determinar $x$ para que o volume do paralelepípedo que tem um dos vértices no ponto $A=(2,1,6)$ e os três vértices adjacentes nos pontos $B=(4,1,3)$, $C=(1,3,2)$ e $D=(1,x,1)$ seja igual a $15$.


%%%%%%%%%%%%%%%%%%%%%%%%%%%%%%%%%%%%%%%%%%%%%%%%%%%%%%%%%%%%%%%%%
\begin{pycode} 
import sympy as sp

v1=sp.Matrix([2,1,-2,1])
v2=sp.Matrix([1,-3,2,-1])
v3=sp.Matrix([3,1,1,0])

b1=3*v1-v2+2*v3
b2=sp.Matrix([1,-1,0,0])

def aumk(X,k):
 x=X.shape[1]
 col = 'c'*(x-k)+'|'+'c'*k
 MM = sp.latex(X).replace(r'\begin{matrix}', r'\begin{array}{'+col+'}')
 MM = MM.replace(r'\end{matrix}', r'\end{array}')
 return MM

A=sp.Matrix.hstack(v1,v2,v3)
B=sp.Matrix.hstack(b1,b2)

M=sp.Matrix.hstack(A,B)
Mb=aumk(M,2)
R=M.rref()[0]
Rb=aumk(R,2)


\end{pycode} 

\question Escreva, se possível,  os vetores {\(\vec{b}_1=\pyl{b1}\)} e {\(\vec{b}_2=\pyl{b2}\)} como combinação linear dos seguintes vetores
\[
\vec{v}_1=\pyl{v1},\ \vec{v}_2=\pyl{v2} \text{ e } \vec{v}_3=\pyl{v3}.
\]

\begin{solution}
Para isso, basta escalonar a matriz aumentada com as duas últimas colunas sendo os vetores \(\vec{b}_1\) e \(\vec{b}_2\).
\begin{equation*}
\py{Mb}\leftrightarrow \py{Rb}
\end{equation*}
Com isso, temos que \(\vec{b}_2\) não pode ser escrito como combinação linear dos vetores dados, uma vez que o sistema considerando-se a matriz dos coeficientes e a última coluna não tem solução. Já o sistema considerando-se a matriz dos coeficientes e a penúltima  nos dá a seguinte solução:
\[\vec{b}_1=3\vec{v}_1-\vec{v}_2+2\vec{v}_3.\]

\end{solution}

%%%%%%%%%%%%%%%%%%%%%%%%%%%%%%%%%%%%%%%%%%%%%%%%%%%%%%%%%%%%%%%%%%%%%%

\begin{pycode}
import sympy as sp

def pt(Y):
 return((Y[0],Y[1],Y[2]))

v1=sp.Matrix([1,1,2])
v2=sp.Matrix([1,0,1])
v3=sp.Matrix([4,6,12])

A=sp.Matrix.hstack(v1,v2,v3)


u1=sp.Matrix([0,3,1,-1])
u2=sp.Matrix([6,0,5,1])
u3=sp.Matrix([4,-7,1,3])

B=sp.Matrix.hstack(u1,u2,u3)
\end{pycode}

\question Para cada caso, diga se o conjunto de vetores é LI ou LD.

\begin{parts}
\part $\{\pyl{pt(v1)},\pyl{pt(v2)}, \pyl{pt(v3)}\}$
\part $\{\pyl{pt(u1)},\pyl{pt(u2)}, \pyl{pt(u3)}\}$
\end{parts}

\begin{solution}
\begin{parts}
\part Basta escalonar a matriz $A$ cujas colunas são so vetores $\vec{v}_1=\pyl{pt(v1)},$ $\vec{v}_2=\pyl{pt(v2)}$ e $\vec{v}_3=\pyl{pt(v3)}$.
\[A=\pyl{A}\leftrightarrow \pyl{A.rref()[0]}\]
Como este sistema tem única solução, os vetores são LI.

\part Basta escalonar a matriz $B$ cujas colunas são os vetores $\vec{v}_1=\pyl{pt(u1)},$ $\vec{v}_2=\pyl{pt(u2)}$ e $\vec{v}_3=\pyl{pt(u3)}.$
\[B=\pyl{B}\leftrightarrow \pyl{B.rref()[0]}\]
Como este sistema tem infinitas solução, os vetores são LD.
\end{parts}
\end{solution}


\question Encontre um subconjunto dos vetores 
\[\vec{v}_1=(1,-2,0,3), \vec{v}_2=(2,-5,-3,6),
\]
\[\vec{v}_3=(0,1,3,0), \vec{v}_4=(2,-1,4,-7), \vec{v}_5=(5,-8,1,2)
\]
que forma uma base para o espaço gerado por eles.

%-------------------------------------------------------

\begin{pycode}
import sympy as sp

A=sp.Matrix([[0,0,2,0],[1,0,1,0],[0,1,-2,0],[0,0,0,1]])
I=sp.eye(4)

l=A.eigenvals()
\end{pycode}

\question Considere a matriz
\[A=
\pyl{A}
\]
\begin{parts}
\part Determine os autovalores de $A$.
\part Determine os autoespaços de $A$.
\part Determine uma base para cada autoespaço e a dimensão.
\end{parts}


\begin{solution} 
\begin{pycode}
import sympy as sp
A = sp.Matrix([
    [0, 0, 2, 0],
    [1, 0, 1, 0],
    [0, 1, -2, 0],
    [0, 0, 0, 1]
])
eigenvals = A.eigenvals()
eigenvects = A.eigenvects()
l=sp.symbols('lambda')

print('Vamos calular o polinômio característico.')
print()

d=sp.factor((A-l*sp.eye(4)).det())
print('\(p(\lambda)=\det(A-\lambda I)=\)',f'\(\det {sp.latex(A-l*sp.eye(4))}={sp.latex(d)}.\)')
print()

print('A seguir são dados os autovalores e a respectiva base do autoespaço.')
print()
for vect in eigenvects:
    print(f"$\lambda={sp.latex(vect[0])}$, autovetores:")
    for v in vect[2]:
        print(f"${sp.latex(v)}$,")
    print()

\end{pycode}


\end{solution}


%-------------------------------------------------------
\begin{pycode}
import sympy as sp

A=sp.Matrix([[4,2,2],[2,4,2],[2,2,4]])
B=sp.Matrix([[-1,7,-1],[0,1,0],[0,15,-2]])

l=sp.symbols('lambda',real=True)
I3=sp.eye(3)
pa=sp.factor((A-l*sp.eye(3)).det())
autvaA = A.eigenvals()
autveA = A.eigenvects()
l1=autveA[0][0]
l2=autveA[1][0]

A1=A-l1*I3
RA1=A1.rref()[0]
sola1=sp.linsolve(RA1,(x,y,z))
sa1 = next(iter(sola1))
vla1 = tuple(sa1.free_symbols)

va1=autveA[0][2][0]
va2=autveA[0][2][1]

A2=A-l2*I3
RA2=A2.rref()[0]
sola2=sp.linsolve(RA2,(x,y,z))
sa2 = next(iter(sola2))
va3=autveA[1][2][0]

Pa=sp.Matrix.hstack(va1,va2,va3)
Da=sp.diag(l1,l1,l2)

\end{pycode}

\question 
\begin{parts}
\part Mostre que 
\[A=\pyl{A}\]
é diagonalizável e obtenha as matrizes diagonal $D$ e a matriz $P$ que diagonaliza $A$.

\part Calcule $B^{11}$, onde 
\[B=\pyl{B}\]
\end{parts}

\begin{solution}
\begin{parts}
\part Primeiramente, vamos calcular o polinômio característico.
\[p(\lambda)=\det(A-\lambda I)=\det\pyl{A-l*sp.eye(3)}=\pyl{pa}.\]
Com isso, temos que os autovalores são \(\lambda_1=\pyl{autveA[0][0]}\) e \(\lambda_2=\pyl{autveA[1][0]}\).
Agora vamos determinar os autoespaços correspondentes.

\noindent\underline{Autoespaço \(E_{\pyl{l1}}=\ker(A-\pyl{l1}I)\):}

\[\pyl{A1}\leftrightarrow \pyl{RA1}\Longrightarrow 
E_{\pyl{l1}}=\left\{ (-\alpha-\beta,\alpha,\beta);\ \alpha,\beta\in \R. \right\}\]
Com isso, \(\dim E_{\pyl{l1}}=2\) e uma base para esse autoespaço é:
\[
\mathcal{B}_{\pyl{l1}}=\left\{ \pyl{va1},\pyl{va2} \right\}
\]

\noindent\underline{Autoespaço \(E_{\pyl{l2}}=\ker(A-\pyl{l2}I)\):}

\[\pyl{A2}\leftrightarrow \pyl{RA2}\Longrightarrow 
E_{\pyl{l2}}=\left\{ (\alpha,\alpha,\alpha);\ \alpha\in \R. \right\}\]
Com isso, \(\dim E_{\pyl{l2}}=1\) e uma base para esse autoespaço é:
\[
\mathcal{B}_{\pyl{l2}}=\left\{ \pyl{va3}\right\}
\]
Portanto, as matrizes $D$ e $P$ são
\[
D=\pyl{Da},\ \ P=\pyl{Pa}.
\]

\end{parts}


\end{solution}



%-------------------------------------------------------------